\documentclass[12pt]{article}
\usepackage[T1]{fontenc}
\usepackage[utf8]{inputenc}
\usepackage{lmodern}
\usepackage{textcomp}
\usepackage{enumitem}
\usepackage{geometry}
\geometry{margin=1in}
\begin{document}
\begin{center}
Answer Key - Health Sciences - Graduate Level Assessment 3 \
\small (Answers are ordered exactly as in the question paper.)
\end{center}

\section*{Multiple Choice}
\begin{enumerate}[label=\arabic*.]
\item \textbf{Answer:} A
\\ \textit{Options:} A) Nitrogen 14; B) Oxygen 16; C) Oxygen 18; D) Helium 3; E) Lithium 6
\\ \textit{Note:} Nitrogen-14 is not actually the correct answer for measuring total body water; the stable isotope commonly used for this purpose is deuterium (a hydrogen isotope), as it can accurately reflect total body water without causing any significant biological effects.
\item \textbf{Answer:} C
\\ \textit{Options:} A) The hypothalamus; B) The parietal lobe; C) The cerebellum; D) The occipital lobe; E) The amygdala
\\ \textit{Note:} The cerebellum is supplied by the posterior inferior cerebellar artery, a branch of the vertebral artery, which itself is a branch of the subclavian artery, thereby establishing the connection between the subclavian arteries and the cerebellum.
\item \textbf{Answer:} A
\\ \textit{Options:} A) Potassium; B) Manganese; C) Iron; D) Zinc; E) Magnesium
\\ \textit{Note:} Potassium is not significantly affected by high intakes of foods containing phytic acid because phytic acid primarily impairs the absorption of minerals such as iron, zinc, and calcium, while potassium absorption remains largely unaffected.
\item \textbf{Answer:} A
\\ \textit{Options:} A) the articular disc of the TMJ and the lateral pterygoid muscle.; B) the capsule and ligaments of the TMJ and the masseter muscle.; C) the ligaments of the TMJ and the medial pterygoid muscle.; D) the articular disc of the TMJ and the medial pterygoid muscle.; E) the articular disc of the TMJ and the medial and lateral pterygoid muscles.
\\ \textit{Note:} The articular disc of the temporomandibular joint (TMJ) and the lateral pterygoid muscle are rich in proprioceptive receptors that provide critical sensory feedback regarding the position and movement of the jaw.
\item \textbf{Answer:} A
\\ \textit{Options:} A) BMI 22-26 kg/m 2; B) BMI 23-27 kg/m 2; C) BMI 25-30 kg/m 2; D) BMI \textgreater{} 30 kg/m 2; E) BMI 24-28 kg/m 2
\\ \textit{Note:} The correct answer defines overweight as a Body Mass Index (BMI) range of 22-26 kg/m², which aligns with health guidelines that categorize individuals within this range as having an excess body weight that may pose health risks.
\end{enumerate}

\section*{True / False}
\begin{enumerate}[label=\arabic*.]
\setcounter{enumi}{5}
\item \textbf{Answer:} False
\\ \textit{Options:} A) True; B) False
\\ \textit{Note:} The muscles of the soft palate are primarily innervated by the vagus nerve (cranial nerve X) and the accessory nerve (cranial nerve XI), rather than the trigeminal and hypoglossal nerves.
\item \textbf{Answer:} False
\\ \textit{Options:} A) True; B) False
\\ \textit{Note:} Brain growth during adolescence is characterized by significant developmental changes and accelerated growth, particularly in the prefrontal cortex and limbic system, which makes it a period of rapid brain maturation rather than slower growth.
\item \textbf{Answer:} True
\\ \textit{Options:} A) True; B) False
\\ \textit{Note:} Feeling unprepared and lacking confidence in the ability to transition successfully into retirement can lead to increased stress, anxiety, and negative health outcomes, making it advisable to delay retirement until one feels more ready.
\item \textbf{Answer:} True
\\ \textit{Options:} A) True; B) False
\\ \textit{Note:} The statement is true because the levator veli palatini muscle elevates the soft palate during swallowing or yawning, which in turn opens the auditory tube to equalize pressure in the middle ear.
\item \textbf{Answer:} True
\\ \textit{Options:} A) True; B) False
\\ \textit{Note:} Diuretic medications promote increased urine production by inhibiting the reabsorption of sodium and water in the kidneys, leading to more frequent urination, which is particularly notable in older adults due to age-related changes in renal function and fluid balance.
\end{enumerate}

\section*{Long Form Response}
\begin{enumerate}[label=\arabic*.]
\setcounter{enumi}{10}
\item \textbf{Answer:} The lower rates of alcohol abuse among Pacific Islander-Americans compared to other demographic groups in the United States can be attributed to a combination of cultural, social, and environmental factors. Many Pacific Islander communities emphasize strong familial ties and cultural traditions that promote healthy lifestyles and discourage substance abuse. Additionally, religious beliefs prevalent in these communities often advocate for sobriety and responsible behavior, further reinforcing positive social norms. Access to community support systems and resources tailored to their specific cultural contexts also plays a significant role in mitigating alcohol-related issues. Finally, socioeconomic factors, including lower exposure to environments where alcohol use is normalized, contribute to these lower rates of alcohol abuse among Pacific Islander- Americans.
\\ \textit{Note:} correct because it identifies the interplay of cultural values, strong family support, and religious influences as key factors that collectively discourage alcohol abuse among Pacific Islander-Americans, differentiating them from other demographic groups.
\item \textbf{Answer:} Accessing the airway through the cricothyroid membrane is a critical procedure in emergency situations, particularly when standard intubation fails or is not possible. The appropriate location for incision is inferior to the thyroid cartilage, specifically within the cricothyroid membrane, which lies between the thyroid and cricoid cartilages. This site is chosen due to its superficial position and the absence of major vascular structures, reducing the risk of significant hemorrhage. Additionally, this approach allows for rapid access to the airway in life-threatening scenarios, facilitating immediate ventilation. Understanding these anatomical considerations is crucial for safe and effective emergency airway management.
\\ \textit{Note:} correct because accessing the airway through the cricothyroid membrane is advantageous in emergencies due to its anatomical positioning, which minimizes the risk of damaging major blood vessels while providing a rapid and effective route for airway access.
\end{enumerate}

\vfill
\noindent\textit{End of Answer Key}
\end{document}